% =====================================================================
\chapter{Variables continuas, desigualdades y ley de los grandes números}
\label{tema:6}
% =====================================================================

Muchas magnitudes, como tiempos, longitudes o pesos de una red neuronal, toman valores en un continuo; su distribución se describe con una densidad en lugar de una función de masa. Este tema introduce las variables continuas y tres distribuciones fundamentales (uniforme, exponencial y normal), y termina con las desigualdades de Markov y Chebyshev y la ley débil de los grandes números, que justifica que una media muestral aproxime una esperanza. Se siguen el capítulo~5 de Blitzstein y Hwang~\cite{bh} (§5.1, §5.2, §5.4 y §5.5) y su §10.1--10.2; las secciones §3.1 a §3.3 y §5.1 a §5.2 de Bertsekas y Tsitsiklis~\cite{bt} (clases~8 y~19 de MIT 6.041~\cite{mit6041f10}); los capítulos~5 y~8 de Ross~\cite{ross}; el capítulo~19 de~\cite{mcs} y el capítulo~8 de Grinstead y Snell~\cite{gs}.

\section{Densidad y función de distribución}

\begin{definicion}[Variable continua]
Una variable aleatoria $X$ es \emph{continua} si existe una función $f_X\colon\R\to[0,\infty)$, su \emph{densidad}, tal que para todos $a\le b$
\[
\PP(a\le X\le b)=\int_a^bf_X(x)\,dx.
\]
Necesariamente $\int_{-\infty}^{\infty}f_X(x)\,dx=1$. La \emph{función de distribución} es $F_X(x)=\PP(X\le x)=\int_{-\infty}^{x}f_X(t)\,dt$.
\end{definicion}

Consecuencias inmediatas: $\PP(X=x)=\int_x^xf_X=0$ para todo $x$, de modo que da igual escribir $<$ o $\le$; $F_X$ es continua y no decreciente, y en los puntos donde $f_X$ es continua, $F_X'(x)=f_X(x)$ por el teorema fundamental del cálculo. La densidad \emph{no} es una probabilidad y puede ser mayor que~1: lo que se interpreta es $\PP(x\le X\le x+\delta)\approx f_X(x)\,\delta$ para $\delta$ pequeño.

\begin{definicion}[Esperanza y varianza]
Para $X$ continua, $\E[X]=\int_{-\infty}^{\infty}x\,f_X(x)\,dx$ (si la integral converge absolutamente) y, para una función $g$, $\E[g(X)]=\int g(x)f_X(x)\,dx$. La varianza se define como en el caso discreto, $\Var(X)=\E[(X-\E X)^{2}]=\E[X^{2}]-\E[X]^{2}$, y la linealidad y las propiedades de la \cref{prop:t5-var} siguen siendo válidas.
\end{definicion}

\section{Distribuciones uniforme y exponencial}

\begin{definicion}[Uniforme]
$X\sim\Unif(a,b)$, con $a<b$, si $f_X(x)=\frac{1}{b-a}$ en $[a,b]$ y $0$ fuera. Entonces $F_X(x)=\frac{x-a}{b-a}$ para $x\in[a,b]$: la probabilidad de un subintervalo es proporcional a su longitud.
\end{definicion}

\begin{proposicion}\label{prop:t6-unif}
Si $X\sim\Unif(a,b)$, entonces $\E[X]=\frac{a+b}{2}$ y $\Var(X)=\frac{(b-a)^{2}}{12}$.
\end{proposicion}

\begin{definicion}[Exponencial]
$X\sim\Expo(\lambda)$, con $\lambda>0$, si $f_X(x)=\lambda e^{-\lambda x}$ para $x\ge0$ y $0$ para $x<0$. Integrando, $F_X(x)=1-e^{-\lambda x}$ y $\PP(X>x)=e^{-\lambda x}$ para $x\ge0$. Se cumple $\E[X]=1/\lambda$ y $\Var(X)=1/\lambda^{2}$ (\cref{eje:6-expo}).
\end{definicion}

\begin{proposicion}[Falta de memoria]\label{prop:t6-memoria}
Si $X\sim\Expo(\lambda)$, entonces $\PP(X>s+t\mid X>s)=\PP(X>t)$ para todos $s,t\ge0$.
\end{proposicion}

La exponencial es la versión continua de la geométrica: modela tiempos de espera hasta un suceso que «no envejece», como la llegada de la siguiente petición a un servidor.

\section{Distribución normal y tipificación}

\begin{definicion}[Normal]
$Z\sim\Nor(0,1)$ (\emph{normal estándar}) si su densidad es $\varphi(z)=\frac{1}{\sqrt{2\pi}}e^{-z^{2}/2}$; su función de distribución se denota $\Phi$. Para $\mu\in\R$ y $\sigma>0$, $X\sim\Nor(\mu,\sigma^{2})$ si
\[
f_X(x)=\frac{1}{\sigma\sqrt{2\pi}}\exp\Bigl(-\frac{(x-\mu)^{2}}{2\sigma^{2}}\Bigr).
\]
\end{definicion}

Que $\varphi$ integra~1 equivale a la integral gaussiana $\int_{-\infty}^{\infty}e^{-z^{2}/2}dz=\sqrt{2\pi}$, que se demuestra en el módulo~5 con coordenadas polares. Por la simetría de $\varphi$, $\E[Z]=0$ y $\Phi(-z)=1-\Phi(z)$; integrando por partes se obtiene $\E[Z^{2}]=1$, luego $\Var(Z)=1$. $\Phi$ no tiene expresión con funciones elementales y se tabula (\cref{tab:t6-phi}).

\begin{teorema}[Tipificación]\label{teo:t6-tipificacion}
Si $Z\sim\Nor(0,1)$, entonces $X=\mu+\sigma Z\sim\Nor(\mu,\sigma^{2})$. Recíprocamente, si $X\sim\Nor(\mu,\sigma^{2})$, la variable tipificada $Z=(X-\mu)/\sigma$ es $\Nor(0,1)$ y
\[
\PP(X\le x)=\Phi\Bigl(\frac{x-\mu}{\sigma}\Bigr).
\]
En particular, $\E[X]=\mu$ y $\Var(X)=\sigma^{2}$.
\end{teorema}

\begin{table}[h]
\centering
\begin{tabular}{@{}lcccccccc@{}}
\toprule
$z$ & $0{,}5$ & $1$ & $1{,}5$ & $1{,}645$ & $1{,}96$ & $2$ & $2{,}5$ & $3$ \\
\midrule
$\Phi(z)$ & $0{,}6915$ & $0{,}8413$ & $0{,}9332$ & $0{,}9500$ & $0{,}9750$ & $0{,}9772$ & $0{,}9938$ & $0{,}9987$ \\
\bottomrule
\end{tabular}
\caption{Valores de la función de distribución normal estándar. Para $z<0$, $\Phi(z)=1-\Phi(-z)$.}
\label{tab:t6-phi}
\end{table}

\addtocounter{definicion}{1}%

\section{Desigualdades de Markov y Chebyshev}

\begin{teorema}[Desigualdad de Markov]\label{teo:t6-markov}
Si $X\ge0$ y $a>0$, entonces $\PP(X\ge a)\le\dfrac{\E[X]}{a}$.
\end{teorema}

\begin{teorema}[Desigualdad de Chebyshev]\label{teo:t6-chebyshev}
Si $X$ tiene media $\mu$ y varianza $\sigma^{2}$ finitas, para todo $c>0$,
\[
\PP(\abs{X-\mu}\ge c)\le\frac{\sigma^{2}}{c^{2}}.
\]
\end{teorema}

Con $c=k\sigma$, Chebyshev dice que $\PP(\abs{X-\mu}\ge k\sigma)\le1/k^{2}$ para \emph{cualquier} distribución: a dos desviaciones típicas hay como mucho un $25\,\%$ de la probabilidad. Para la normal el valor exacto es un $4{,}6\,\%$; la cota es peor porque no usa nada de la forma de la distribución.

\section{Ley débil de los grandes números}

\begin{proposicion}[Media muestral]\label{prop:t6-media}
Sean $X_1,\dots,X_n$ independientes con la misma media $\mu$ y la misma varianza $\sigma^{2}$, y sea $\Xbar_n=\frac1n\sum_{i=1}^{n}X_i$. Entonces $\E[\Xbar_n]=\mu$ y $\Var(\Xbar_n)=\sigma^{2}/n$.
\end{proposicion}

\begin{teorema}[Ley débil de los grandes números]\label{teo:t6-lgn}
En las condiciones de la \cref{prop:t6-media}, para todo $\varepsilon>0$,
\[
\PP\bigl(\abs{\Xbar_n-\mu}\ge\varepsilon\bigr)\le\frac{\sigma^{2}}{n\varepsilon^{2}}\xrightarrow[n\to\infty]{}0.
\]
Se dice que $\Xbar_n$ \emph{converge en probabilidad} a $\mu$.
\end{teorema}

\begin{corolario}[Frecuencias relativas]\label{cor:t6-frecuencias}
Si $A$ es un suceso de probabilidad $p$ y $f_n$ es la frecuencia relativa con que ocurre en $n$ repeticiones independientes, entonces $\PP(\abs{f_n-p}\ge\varepsilon)\le\dfrac{p(1-p)}{n\varepsilon^{2}}\le\dfrac{1}{4n\varepsilon^{2}}$.
\end{corolario}

El corolario reconcilia la definición axiomática con la interpretación frecuentista: la probabilidad de un suceso es el valor al que tiende, en probabilidad, su frecuencia relativa. Con $\varepsilon=0{,}01$ y $n=10^{5}$, la cota es $1/(4\cdot10^{5}\cdot10^{-4})=0{,}025$.

\addtocounter{definicion}{1}%

\addtocounter{definicion}{1}%

